\begin{annexitem}{annex:p3.ab.1}{thm:p03-ab-index-normal}{Effective index reduction}
\begin{proof}
Specialize \cref{thm:p03-w03-p03-073-pp-elimination} to
\(\mathbb Z\)-modules. A pp sentence is true in every module: assign
zero to all its quantified variables. Consequently the theorem gives,
for each sentence, a Boolean combination of finite pp-index thresholds
with no free variables, equivalent in every abelian group. Translate
integer scalar notation back into the group language by moving negative
summands across equations. This is an effective definitional translation.

The existence theorem suffices to obtain an algorithm without asking
whether an arbitrary proposed equivalence is valid. Enumerate all finite
Boolean combinations of pp-index thresholds, and dovetail an enumeration
of all finite first-order proofs from the finite axioms \(\mathsf{Ab}\).
Stop when a proof of an equivalence between the input and one enumerated
combination appears. Proof checking and enumeration are effective. The
normal-form theorem and first-order completeness ensure that some such
proof exists, so this search terminates. Replace each pair \((H,J)\) by
\((H,H\land J)\), writing its second member as \(K\).

Initially these pp subgroups may be subgroups of \(G^r\). For \(r>1\),
let \(\pi:G^r\longrightarrow G^{r-1}\) forget the last coordinate.
There is an exact sequence of abelian groups
\[
 0\longrightarrow
 (H\cap\ker\pi)/(K\cap\ker\pi)
 \longrightarrow H/K
 \longrightarrow \pi H/\pi K\longrightarrow0.
\]
The last arrow sends \(h+K\) to \(\pi h+\pi K\). If its value is zero,
choose \(k\in K\) with \(\pi h=\pi k\); then \(h-k\in H\cap\ker\pi\).
This proves exactness, including the asserted kernel. The subgroups
\(H\cap\ker\pi\) and \(K\cap\ker\pi\), identified with subgroups of
the last copy of \(G\), are unary pp subgroups, obtained by substituting
zero for the other coordinates. The projected subgroups are pp definable
by existentially quantifying the last coordinate. Repeat on the latter
pair. In arity zero the index is one.

For a finite threshold \(k\), the exact sequence determines whether
\([H:K]\geq k\) from the indices of these factors capped at \(k\).
Indeed, if either factor has at least \(k\) elements then so does the
middle group. Otherwise both factors are finite and the middle group
has the product of their orders, since all fibers have the kernel's
order. A capped value \(j<k\) is specified by the two thresholds
\(\text{index}\geq j\) and \(\text{index}\not\geq j+1\), with the
first condition omitted for \(j=1\); value \(k\) is specified by
\(\text{index}\geq k\). Enumerating the finitely many capped factor
values whose product reaches \(k\) gives a Boolean combination of unary
index thresholds. Recursion on the arity terminates. Applying these
replacements to the enumerated normal form proves the assertion.
\end{proof}
\end{annexitem}

\begin{annexitem}{annex:p3.ab.2}{thm:p03-ab-unary-normal}{Integer reduction}
\begin{proof}
Write the unary pp formula as
\(\exists\boldsymbol y\,(A x+B\boldsymbol y=0)\), where \(A\) is an
integer column and \(B\) a finite integer matrix. Integer elementary row
operations are invertible changes of the list of equations in every
abelian group. Integer elementary column operations on \(B\) are
invertible changes of the quantified variables. In particular they do
not require division in the group.

For completeness, diagonalization over the integers is effective as
follows. In a nonzero remaining submatrix move a nonzero entry to its
upper left corner and make it positive. Divide each other entry in its
column by this pivot and subtract the quotient times the pivot row.
A nonzero remainder can be moved to the pivot position and strictly
reduces the positive pivot. Apply the corresponding column operations
to the pivot row as well. Each failure of exact divisibility strictly
reduces the pivot; hence after finitely many reductions the pivot divides
all entries in its row and column. Clear that row and column by exact
subtractions and continue on the smaller submatrix. All operations are
unimodular. No divisibility ordering of the final diagonal entries is
needed. This gives computable unimodular matrices \(U,V\) with
\(UBV\) diagonal, with positive nonzero entries \(d_1,\ldots,d_t\).

Put \(\boldsymbol y=V\boldsymbol z\) and multiply the equations by \(U\).
Each nonzero diagonal row has the form \(d_i z_i=c_i x\), with a possible
sign absorbed into \(c_i\). Existentially quantifying its variable gives
\(d_i\mid c_i x\). A row with zero quantified coefficients gives
\(e_jx=0\). The conjunction of these latter equations is equivalent to
\(mx=0\), where \(m\) is the nonnegative greatest common divisor of
the \(e_j\), or zero if all are zero or there are no such rows. One
implication follows because every \(e_j\) is a multiple of \(m\); the
other follows from a finite Bezout identity for their greatest common
divisor. Unused quantified variables can be zero, and a row with
\(c_i=0\) imposes no condition. This proves the claimed uniform form.
\end{proof}
\end{annexitem}

\begin{annexitem}{annex:p3.ab.3}{thm:p03-ab-bounded}{Bounded groups}
\begin{proof}
First suppose \(p^E A=0\), with \(E\geq1\). Choose an
\(\mathbb F_p\)-basis \((b_i)_{i\in I}\) of \(p^{E-1}A\), and choose
\(a_i\in A\) with \(p^{E-1}a_i=b_i\). These elements generate a subgroup
\(B\cong\bigoplus_{i\in I}C_{p^E}\). To verify independence, suppose
there is a nonzero finite relation with coefficients reduced modulo
\(p^E\), and let \(t<E\) be the smallest valuation of its nonzero
coefficients. Multiplying the relation by \(p^{E-1-t}\) gives a nonzero
linear relation over \(\mathbb F_p\) between the \(b_i\), a contradiction.

There is a retraction \(A\longrightarrow B\). Here is the extension
argument, so no splitting theorem is assumed. Order pairs \((C,f)\),
where \(B\subseteq C\subseteq A\) and \(f:C\longrightarrow B\) restricts
to the identity on \(B\), by extension. Unions of chains are again such
pairs, so choose a maximal pair by Zorn's lemma. If \(a\notin C\), the
coset \(a+C\) has order \(p^k\) for some \(1\leq k\leq E\). The element
\(f(p^ka)\) is killed by \(p^{E-k}\). In \(B\), whose cyclic summands
all have order \(p^E\), every element killed by \(p^{E-k}\) is a
\(p^k\)-multiple. Choose \(b\in B\) with \(p^kb=f(p^ka)\). Then
\[
 f'(c+na)=f(c)+nb
\]
defines an extension to \(C+\mathbb Za\): a relation \(na\in C\)
forces \(p^k\mid n\), which proves well-definedness using the choice
of \(b\). This contradicts maximality. Thus \(C=A\).

Write \(A=B\oplus A'\) using the retraction. Since
\(p^{E-1}A=p^{E-1}B\), the element \(p^{E-1}a'\) for \(a'\in A'\)
lies both in \(B\) and in the kernel of the retraction, and is zero.
Induction on \(E\), with the zero group as the case \(E=0\), therefore
expresses \(A\) as a direct sum of cyclic groups of orders
\(p,\ldots,p^E\).

Images under multiplication, kernels of multiplication, and their
intersections commute with direct sums. In a cyclic group of order
\(p^s\), the space \(U_{p,n}/U_{p,n+1}\) has dimension one when
\(n=s-1\) and zero otherwise. It follows that the multiplicity of
\(C_{p^s}\) in the decomposition is precisely \(\alpha_{p,s-1}(A)\),
including when that dimension is infinite.

Finally suppose \(mA=0\) and factor \(m=\prod_p p^{E_p}\). The Chinese
remainder identity gives integer coefficients which act on \(A\) as
pairwise orthogonal projections summing to the identity, with the
\(p\)-projection having image \(A[p^{E_p}]\). Explicitly choose an
integer congruent to one modulo \(p^{E_p}\) and zero modulo each other
prime-power factor; multiplication by it is that projection. Thus
\(A=\bigoplus_{p\mid m}A[p^{E_p}]\). The case \(m=1\) is the zero
group. This completes the bounded-exponent decomposition.
\end{proof}
\end{annexitem}

\begin{annexitem}{annex:p3.ab.4}{thm:p03-ab-layers}{Layer comparisons}
\begin{proof}
Fix \(p\) and abbreviate \(U_{p,n}(G)\) to \(U_n\). The first recurrence
is dimension additivity for the subspace \(U_{n+1}\subseteq U_n\).
For the second, multiplication by \(p\) gives a surjective linear map
\[
 p^nG/p^{n+1}G\longrightarrow p^{n+1}G/p^{n+2}G.
\]
If \(x\in p^nG\) maps to zero, write \(px=p^{n+2}y\). Then
\(x-p^{n+1}y\in U_n\) and represents the same class as \(x\). The
kernel is consequently the image of \(U_n\), whose kernel in this
quotient is \(U_n\cap p^{n+1}G=U_{n+1}\). This gives a short exact
sequence with first term \(U_n/U_{n+1}\). Dimension additivity proves
the second recurrence, as a cardinal sum even for infinite dimensions.

Set \(A=G[p^E]\). For \(0\leq n<E\), one has
\(U_{p,n}(A)=U_n\). Indeed if \(x=p^ny\) and \(px=0\), then
\(p^{n+1}y=0\), so \(y\in A\); the reverse inclusion is immediate.
Furthermore \(p^EA=0\), so the top \(U\)-space after level \(E-1\)
is zero. Apply the bounded-group theorem: below the top the dimensions
are \(\alpha_{p,s-1}(G)\), and at the top they are
\(\dim U_{E-1}=\beta_{p,E-1}(G)\).

Next put \(A=G/p^EG\). For \(0\leq j\leq E-2\), the natural map induces
an isomorphism
\[
 U_j(G)/U_{j+1}(G)\longrightarrow U_{p,j}(A)/U_{p,j+1}(A).
\]
To prove surjectivity, represent an element of \(U_{p,j}(A)\) by
\(x\in p^jG\) with \(px=p^Ez\). Then
\(t=x-p^{E-1}z\) belongs to \(U_j(G)\); the subtracted element represents
an element of \(U_{p,j+1}(A)\), since \(E-1\geq j+1\) and its
\(p\)-multiple is zero modulo \(p^EG\). For the kernel, if
\(t\in U_j(G)\) has image in \(p^{j+1}A\), then
\(t\in p^{j+1}G+p^EG=p^{j+1}G\), hence \(t\in U_{j+1}(G)\).
Conversely every such element maps to \(U_{p,j+1}(A)\). This proves
the isomorphism. At the top,
\[
 U_{p,E-1}(A)=p^{E-1}A
             \cong p^{E-1}G/p^EG,
\]
whose dimension is \(\gamma_{p,E-1}(G)\). The bounded-group theorem
now gives the asserted quotient row, also when \(E=1\), for which
there are no lower levels to check.
\end{proof}
\end{annexitem}

\begin{annexitem}{annex:p3.ab.5}{thm:p03-ab-index-calculation}{Index calculations}
\begin{proof}
Use the unary forms supplied by \cref{thm:p03-ab-unary-normal} and
replace \(K\) by \(H\land K\). Thus inclusion is maintained without
any decision of implication between formulas. Write \(m_H,m_K\) for
the annihilation coefficients of the resulting forms. If \(m_H>0\),
conjoining its equation into \(K\) ensures that the latter also has a
positive annihilation coefficient (a greatest common divisor including
\(m_H\)). Let \(D\) be as in the stated algorithm.

First verify the cyclic calculation. In \(C_{p^s}\), for \(m>0\),
the equation \(mx=0\) requires
\[
 x\in p^{\max(s-v_p(m),0)}C_{p^s}.
\]
For \(m=0\) it imposes no restriction. Multiplication by \(d_i\) has image
\[
 p^{\min(s,v_p(d_i))}C_{p^s}.
\]
Consequently \(d_i\mid c_ix\) requires \(x\) to have valuation at least
\[
 \min(s,v_p(d_i))-v_p(c_i),
\]
with negative requirements vacuous. Intersecting the resulting subgroups
gives exactly \(h_\varphi(p,s)\), which lies between zero and \(s\).
The inclusion \(K\subseteq H\) gives a nonnegative weight \(w=h_K-h_H\)
and a cyclic quotient of order \(p^w\).

A pp solution subgroup commutes with arbitrary direct sums. One
inclusion follows by projecting a solution and its finitely many
witnesses to each coordinate. For the other, an element of a direct
sum has finite support. Coordinate witnesses in that finite support
can be assembled into direct-sum witnesses, taking zero elsewhere;
all the equations are homogeneous. Consequently the quotient \(H/K\)
on a direct sum is the direct sum of the coordinate quotients.
For a bounded group with cyclic multiplicities \(\kappa_{p,s}\), a
finite threshold \(k\) is therefore determined by the capped values
\(\operatorname{cap}_k(\kappa_{p,s})\) and the stated finite product.
If every positive-weight multiplicity is below \(k\), all contributing
quotients are finite and the product is their exact order. If some
positive-weight multiplicity is at least \(k\), it supplies at least
\(k\) nontrivial cyclic summands and in particular at least \(k\)
elements; the capped product also reaches \(k\). Zero weights contribute
nothing, regardless of the corresponding cardinal. This reasoning does
not identify the cardinality of an infinite direct sum with a finite-base
cardinal exponent.

Suppose first that \(m_H=m_K=0\). Every divisibility condition occurring
in either formula is unchanged on adding an element of \(DG\), and it
holds for every element of \(DG\). Moreover it is detected in \(G/DG\).
Indeed if \(c_ix=d_iy+D z\), then \(d_i\mid D\) lets the last summand
be absorbed into \(d_iy\). It follows that \(DG\subseteq K\subseteq H\)
and that their images in \(G/DG\) are exactly the subgroups defined by
the same formulas there. Thus the required index is unchanged. By the
Chinese remainder map,
\[
 G/DG\cong\bigoplus_{p^E\parallel D}G/p^EG.
\]
This isomorphism follows from the usual integer Bezout projections; it
is valid for every abelian group, not merely free groups. The quotient
row of \cref{thm:p03-ab-layers} supplies all the cyclic multiplicities.

Suppose next that \(m_H>0\), and put \(L=m_HD\). Every element of
\(H(G)\), and hence of \(K(G)\), lies in \(G[L]\). If such an element
satisfies \(d_i y=c_i x\) in \(G\), then
\(m_Hd_i y=c_i m_Hx=0\), so \(y\in G[L]\). Therefore every necessary
witness already lies in \(G[L]\), and the same unary formulas define
exactly \(H(G)\) and \(K(G)\) within this bounded subgroup. Decompose
\(G[L]\) into its primary torsion subgroups and use the torsion row
of the layer theorem.

Finally suppose \(m_H=0<m_K\), let \(k\geq2\), and put
\(L=m_KD(k-1)!\). If \([H:K]<k\), the quotient is a finite group of
order at most \(k-1\), and hence is killed by \((k-1)!\). Thus
\((k-1)!H\subseteq K\), and \(m_K(k-1)!H=0\). Since \(DG\subseteq H\),
this implies \(LG=0\). Contrapositively the index threshold is true
whenever \(\neg\operatorname{Ann}_L\) holds. On the remaining branch,
\(G\) is bounded by \(L\), so its multiplicities are exactly the
\(\alpha_{p,s-1}(G)\) in the bounded-group theorem. The cyclic product
calculation applies directly to \(G\).

All prime factorizations and valuation calculations concern supplied
finite integers and terminate. To express a capped dimension value
\(j<k\), use \((\dim\geq j)\land\neg(\dim\geq j+1)\), omitting the
first conjunct for \(j=0\); express value \(k\) by \(\dim\geq k\).
Enumerating all capped assignments passing the product test gives a
finite disjunction of such conditions. Include the appropriate
annihilator branch in the third case. The threshold \(k=1\) is always
true, since a subgroup quotient is nonempty. This establishes both
uniform equivalence and an effective construction in every case.
\end{proof}
\end{annexitem}

\begin{annexitem}{annex:p3.ab.6}{thm:p03-ab-realization}{Finite realizability}
\begin{proof}
It suffices to test a conjunction of literals after a finite Boolean
conversion to disjunctive normal form. Take the bound \(b\) in the
algorithm. Threshold zero is always true and can be removed before
this conversion. For cardinals \(\kappa_1,\ldots,\kappa_t\) and a
finite \(r\leq b\),
\[
 \kappa_1+\cdots+\kappa_t\geq r
 \quad\Longleftrightarrow\quad
 \operatorname{cap}_b(\kappa_1)+\cdots+
 \operatorname{cap}_b(\kappa_t)\geq r.
\]
Indeed if any summand is at least \(b\), both sides hold; otherwise the
sums are equal finite sums. This observation preserves both positive
and negative dimension-threshold literals.

If there are positive annihilator literals, their conjunction is
\(dG=0\), where \(d\) is their greatest common divisor: use Bezout
for one direction and divisibility for the other. Every possible model
therefore has a bounded-group decomposition with cyclic orders dividing
\(d\). On such a decomposition,
\[
 \alpha_{p,n}=\kappa_{p,n+1},\qquad
 \beta_{p,n}=\gamma_{p,n}=\sum_{s>n}\kappa_{p,s},
\]
with absent multiplicities zero. Replace each multiplicity by its cap
at \(b\). This preserves all measured dimension literals by the preceding
observation. It also preserves every annihilator literal: a direct sum
of these cyclic groups is killed by \(m\) exactly when each cyclic
order having nonzero multiplicity divides \(m\). Capping at \(b\geq1\)
preserves whether a multiplicity is zero. Thus every satisfying group
has a satisfying representative in the first finite family. Conversely
any candidate passing the tests is an actual satisfying group. This
also covers \(d=1\), when the family consists of the trivial group.

Now suppose that there is no positive annihilator literal. For each
measured prime \(p\), let \(N_p\) be the largest level appearing in a
dimension literal. Direct calculation gives the following data for
the groups in the second family. In \(C_{p^{j+1}}\), \(\alpha_{p,n}\)
is one exactly at \(n=j\), and \(\beta_{p,n}\) and \(\gamma_{p,n}\)
are one exactly for \(n\leq j\). In \(C_{p^\infty}\), multiplication
by \(p\) is surjective and the subgroup killed by \(p\) has order
\(p\); hence every \(\beta_{p,n}\) is one and every \(\alpha_{p,n}\)
and \(\gamma_{p,n}\) is zero. The group \(\mathbb Z_{(p)}\) is
torsion-free and
\(p^n\mathbb Z_{(p)}/p^{n+1}\mathbb Z_{(p)}\cong\mathbb F_p\),
so its \(\gamma_{p,n}\) is one and its \(\alpha,\beta\) values are
zero. For a different prime \(q\), multiplication by \(q\) is
bijective on each of these \(p\)-indexed groups, so all their
\(q\)-data are zero. On \(\mathbb Q\) all three kinds of dimension
are zero for every prime. Images and kernels commute with these direct
sums, so their data add. This proves the formulas in the algorithm.

If an arbitrary group \(G\) satisfies the conjunction, choose
\[
 \begin{aligned}
 a_{p,j}&=\operatorname{cap}_b(\alpha_{p,j}(G))
                    &&(0\leq j\leq N_p),\\
 b_p&=\operatorname{cap}_b(\beta_{p,N_p+1}(G)),&
 c_p&=\operatorname{cap}_b(\gamma_{p,N_p+1}(G)).
 \end{aligned}
\]
Iterating the two layer recurrences through the finite interval from
\(n\) to \(N_p\) gives
\[
 \begin{aligned}
 \beta_{p,n}(G)&=\sum_{j=n}^{N_p}\alpha_{p,j}(G)
                         +\beta_{p,N_p+1}(G),\\
 \gamma_{p,n}(G)&=\sum_{j=n}^{N_p}\alpha_{p,j}(G)
                         +\gamma_{p,N_p+1}(G).
 \end{aligned}
\]
The capped representative therefore has exactly the same truth values
for all measured dimension literals. Its rational summand ensures
\(mG\ne0\) for every positive integer \(m\), and consequently satisfies
all negative annihilator literals. There are no positive ones in this
case. This constructs a representative in the finite family whenever
a model exists. Notice that the argument only matches finitely many
thresholds. It neither classifies arbitrary groups by an asserted direct
sum decomposition nor subtracts infinite cardinals nor takes a limit
of their heights.

There are finitely many primes, levels and multiplicity assignments
in either family. Their invariants and annihilator truth values are
given by the displayed finite integer calculations, so each candidate
can be checked effectively. The families use only finitely many
countable summands with finite multiplicities and have finite
specifications. Testing all candidates for each conjunction terminates
and returns a countable group description exactly in the satisfiable
case.
\end{proof}
\end{annexitem}

\begin{annexitem}{annex:p3.ab.7}{thm:p03-abelian-decidable}{The decision procedure}
\begin{proof}
Apply effective index reduction to the input sentence, then integer
reduction and the three index calculations to its unary thresholds.
Each algorithm terminates and preserves truth in every abelian group.
Apply finite realizability to the resulting Boolean combination.
Its answer decides satisfiability. To decide validity, apply the same
procedure to the negation of the input: validity holds exactly when
that negation has no abelian-group model. If validity fails, the latter
run returns the stated countable countermodel description. First-order
completeness identifies validity in all models of \(\mathsf{Ab}\) with
provability from those axioms, so this also decides their deductively
closed theory. No decision procedure for the invariant sequence of a
specified arbitrary group is used at any stage.
\end{proof}
\end{annexitem}

\begin{annexitem}{annex:p3.ab.8}{thm:p03-ab-individual-boundary}{An individual theory}
\begin{proof}
Choose a nonrecursive subset \(A\) of the primes, using their computable
enumeration to identify subsets with subsets of \(\mathbb N\), and let
\(G=\bigoplus_{p\in A}C_p\). This is a countable abelian group. For any
prime \(p\), the effectively constructed sentence
\[
 \exists x\,(px=0\land x\ne0)
\]
holds in \(G\) exactly when \(p\in A\). If \(p\notin A\), multiplication
by \(p\) is injective on every summand and hence on their direct sum;
if \(p\in A\), a generator in that summand witnesses the sentence.
A decision procedure for the complete theory of this particular group
would therefore decide \(A\), a contradiction.
\end{proof}
\end{annexitem}
